El proyecto en desarrollo presenta tanto una metodología matemática nueva en el área de predicción fenológica como resultados con potencial impacto en el campo agrícola, dado que se trata de un calendario agrícola basado en un modelo de autoría propia. Por tanto, para la publicación del trabajo como un artículo de investigación en una revista científica se han considerado aquellas revistas cuyos alcances, objetivos y publicaciones previas sean compatibles con el balance entre un desarrollo metodológico teórico y un análisis bajo condiciones reales.\\

Con lo anterior y con base en la literatura consultada hasta el momento, en la siguiente tabla se presentan las revistas indexadas propuestas para someter nuestro trabajo en el mediano plazo.\\


\begin{table}[h]
	\centering
	\begin{tabular}{p{7.5cm}|c|c}
		\hline
		\textbf{Revista} & \textbf{Factor de impacto} & \textbf{Referencias consultadas} \\
		\hline
		\href{https://www.sciencedirect.com/journal/computers-and-electronics-in-agriculture/about/aims-and-scope}{Computers and Electronics in Agriculture}\newline
		Editorial: Elsevier B.V.
		& 10.3
		& \cite{Elnesr_2013} \\ \hline
		\href{https://www.sciencedirect.com/journal/sustainable-computing-informatics-and-systems}{Sustainable Computing: Informatics and Systems}\newline
		Editorial: Elsevier B.V.
		& 6.2
		& \cite{GUMUSCU2020100308} \\ \hline
		\href{https://onlinelibrary.wiley.com/journal/1439037x}{Journal of Agronomy and Crop Science}\newline
		Editorial: Wiley.
		& 3.2
		& \cite{Parker_2016} \\ 
		\hline
	\end{tabular}
	\caption{Revistas idexadas en el Journal Citation Reports (JCR)   consideradas para la publicación del artículo, factor de impacto y las referencias relacionadas con el proyecto consultadas en ellas.}
	\label{tab:journals}
\end{table}


Derivado del avance obtenido y considerando las actividades pendientes de programación, redacción y revisión, se propone diciembre de 2027 como fecha tentativa para el envío del manuscrito a una de las revistas antes mencionadas, correspondiendo con el final del quinto semestre del doctorado.
